\begin{figure*}[t]
\begin{tcolorbox}[colback=white, width=\linewidth, breakable, enhanced, 
                  title=Prompt for the GUI-Owl-1.5 Model, fonttitle=\bfseries]

\lstset{
basicstyle=\scriptsize\ttfamily,
    columns=fullflexible,
    breaklines=true,
    breakatwhitespace=false,
    showstringspaces=false,
    lineskip=-1pt
}

\begin{lstlisting}
# Tools

You may call one or more functions to assist with the user query.

You are provided with function signatures within <tools></tools> XML tags:
<tools>
{"type": "function", "function": {"name": "computer_use", "description": "Use a mouse to interact with a computer.\n* The screen's resolution is 1000x1000.\n* Make sure to click any buttons, links, icons, etc with the cursor tip in the center of the element. Don't click boxes on their edges unless asked.", "parameters": {"properties": {"action": {"description": "The action to perform. The available actions are:\n* `left_click`: Click the left mouse button with coordinate (x, y) pixel coordinate on the screen.", "enum": ["left_click"], "type": "string"}, "coordinate": {"description": "(x, y): The x (pixels from the left edge) and y (pixels from the top edge) coordinates to click.", "type": "array"}}, "required": ["action"], "type": "object"}}}
</tools>

For each function call, return a json object with function name and arguments within <tool_call></tool_call> XML tags:
<tool_call>
{"name": <function-name>, "arguments": <args-json-object>}
</tool_call>

## User Instruction
{instruction}

\end{lstlisting}

\end{tcolorbox}
\end{figure*}